%% preamble.tex — packages and macros for the reviewer response letter
%% Only what a response document actually needs; heavier paper macros are omitted.
%% ─────────────────────────────────────────────────────────────────────────────


%% ── Core packages ────────────────────────────────────────────────────────────
\usepackage[utf8]{inputenc}
\usepackage{xspace}
\usepackage{hyperref}
\usepackage{color}


%% ── To-do / author notes ─────────────────────────────────────────────────────
%% Swap the two lines to hide all notes in the submitted version.
\usepackage[]{todonotes}
% \usepackage[disable]{todonotes}

% Author notes — rename / add as needed
\newcommand{\authorA}[1]{\todo[inline,color=green!40]{\footnotesize{AUTHOR A: #1}}}
\newcommand{\authorB}[1]{\todo[inline,color=red!40]{\footnotesize{AUTHOR B: #1}}}
\newcommand{\authorC}[1]{\todo[inline,color=blue!40]{\footnotesize{AUTHOR C: #1}}}


%% ── Core response macros ─────────────────────────────────────────────────────

% Typeset a reviewer comment in dark gray
\newcommand{\review}[1]{%
    \par\noindent{\color{darkgray}{\textbf{Reviewer comment:} #1}}\smallskip}

% Typeset the authors' response in blue
\newcommand{\reply}[1]{%
    \par\noindent{\color{blue}{\textbf{Authors' response:} #1}}\bigskip}

% Inline reviewer note (for internal use during drafting)
\newcommand{\reviewer}[2]{%
    \todo[color=red!40, inline]{\footnotesize{\##1: #2}}}


%% ── Change-tracking ──────────────────────────────────────────────────────────
%% \changed{} marks text revised in the manuscript.
%% Switch to the colour variant when building the revision to make diffs visible.
\newcommand{\changed}[1]{#1}
% \newcommand{\changed}[1]{{\color{red}#1}}

%% \changedTo{} always renders in red (for text rewritten from scratch).
\newcommand{\changedTo}[1]{{\color{red}#1}}


%% ── Inline quote helper ──────────────────────────────────────────────────────
\newcommand{\say}[1]{``\emph{#1}''}


%% ── Cross-reference shorthands ───────────────────────────────────────────────
\newcommand{\fig}[1]{Fig.~\ref{#1}}
\newcommand{\tab}[1]{Table~\ref{#1}}
\newcommand{\sect}[1]{Section~\ref{#1}}


%% ── Latin abbreviations ──────────────────────────────────────────────────────
\newcommand{\etal}{et al.\xspace}
\newcommand{\ie}{i.e.,\xspace}
\newcommand{\eg}{e.g.,\xspace}


%% ── Domain shorthands (keep in sync with paper/preamble.tex) ─────────────────
%% Add or remove to match your paper's preamble so \entity{} calls work here too.
\newcommand{\entity}[1]{\textsf{#1}\xspace}
